\documentclass[11pt]{article}

\usepackage[a4paper,margin=1in]{geometry}
\usepackage[T1]{fontenc}
\usepackage[utf8]{inputenc}
\usepackage{lmodern}
\usepackage{microtype}
\usepackage{amsmath,amssymb,amsthm,mathtools}
\usepackage{xurl}
\usepackage{authblk}
\usepackage{xcolor}
\usepackage[
  colorlinks=true,
  linkcolor=blue!55!black,
  citecolor=blue!55!black,
  urlcolor=blue!55!black,
  unicode=true,
  pdftitle={A Proposed Solution to Erdős Problem 486},
  pdfauthor={Shouqiao Wang}
]{hyperref}

\numberwithin{equation}{section}

\newtheorem{theorem}{Theorem}[section]
\newtheorem{proposition}[theorem]{Proposition}
\newtheorem{lemma}[theorem]{Lemma}
\newtheorem{corollary}[theorem]{Corollary}
\theoremstyle{definition}
\newtheorem{definition}[theorem]{Definition}
\theoremstyle{remark}
\newtheorem{remark}[theorem]{Remark}

\newcommand{\N}{\mathbb N}
\newcommand{\Z}{\mathbb Z}
\newcommand{\Pp}{\mathbb P}
\newcommand{\E}{\mathbb E}
\newcommand{\Zh}{\widehat{\mathbb Z}}
\newcommand{\cS}{\mathcal S}
\newcommand{\cI}{\mathcal I}

\title{A Proposed Solution to Erd\H{o}s Problem 486}
\author{Shouqiao Wang}
\affil{Columbia University\hspace{2.5em}Multiscalar Intelligence}
\date{}

\begin{document}
\maketitle

\begin{abstract}
For \(A\subseteq\N\) and arbitrary sets
\(X_n\subseteq\Z/n\Z\), let
\[
B=\{m\in\N:m\bmod n\notin X_n
       \text{ for every }n\in A\text{ with }n<m\}.
\]
Erd\H{o}s asked whether \(B\) must possess a logarithmic density.  We
propose a negative answer.  The main ingredient is a finite probabilistic
construction: at each sufficiently large dyadic scale, a fixed positive
proportion of the integers in a short interval is deleted by moduli that are already active,
whereas the union of the associated periodic residue cylinders has
stretched-exponentially small density.  A gliding-hump construction then
places long runs of such blocks between recovery gaps and produces one
fixed congruence system for which
\[
\liminf_{x\to\infty}\frac1{\log x}
  \sum_{\substack{m<x\\m\in B}}\frac1m
\le\frac{177}{200}<\frac{49}{50}\le
\limsup_{x\to\infty}\frac1{\log x}
  \sum_{\substack{m<x\\m\in B}}\frac1m .
\]
This proposed solution was found by GPT-5.6.
\end{abstract}

\section{Introduction}

Throughout, \(\N=\{1,2,\ldots\}\), and all logarithms are natural unless a
base is explicitly indicated.  Erd\H{o}s stated the question, with
the activation condition built into the formulation, in
\cite[p.~48]{Er57}; he later recorded it as Problem~I.26 in
\cite[pp.~235--236]{Er61}.
For each \(n\) in a set \(A\subseteq\N\), choose an arbitrary set
\(X_n\subseteq\Z/n\Z\), and define
\begin{equation}\label{eq:defB}
B=B(A,X):=\{m\in\N:m\bmod n\notin X_n
      \text{ for every }n\in A\text{ with }n<m\}.
\end{equation}
The problem asks whether the logarithmic averages
\[
L_B(x):=\frac1{\log x}
   \sum_{\substack{m<x\\m\in B}}\frac1m
\]
must converge as \(x\to\infty\).

\begin{theorem}[Main theorem]\label{thm:main}
There are a fixed infinite \(A\subseteq\N\) and fixed
\(X_n\subseteq\Z/n\Z\), \(n\in A\), for which the set \(B\) in
\eqref{eq:defB} has no logarithmic density.  More precisely,
\[
\liminf_{x\to\infty}L_B(x)\le\frac{177}{200}
\quad\text{and}\quad
\limsup_{x\to\infty}L_B(x)\ge\frac{49}{50}.
\]
\end{theorem}

\begin{remark}[Strict and inclusive activation]
Some sources use \(n\le m\) in place of \(n<m\).  Let \(B_{\ge}\) and
\(B_{>}\) denote the corresponding survivor sets.  Plainly
\(B_{\ge}\subseteq B_{>}\).  If
\(m\in B_{>}\setminus B_{\ge}\), then \(m\in A\) and \(0\in X_m\).
Moreover, \(B_{>}\setminus B_{\ge}\) is primitive: if \(d<e\) were two of
its elements and \(d\mid e\), then the active modulus \(d<e\), together
with \(0\in X_d\), would exclude \(e\) from \(B_{>}\).
Behrend's theorem gives
\[
 \sum_{\substack{d<x\\d\in B_{>}\setminus B_{\ge}}}\frac1d=o(\log x)
\]
for every primitive set \cite{Be35}.  Thus the two conventions are
equivalent for the existence and value of logarithmic density.
\end{remark}

The zero-residue case \(X_n=\{0\}\) has an affirmative answer by the
Davenport--Erd\H{o}s theorem \cite{DaEr36}; a later elementary proof was
given in \cite{DaEr51}.  Besicovitch showed that natural density may fail
to exist in that setting \cite{Be34}.  More recently, Ara\'ujo proved, in
the inclusive activation convention, that a delayed multi-residue sieve has
a natural density under the summability hypothesis
\[
 \sum_{n\in A}\frac{|X_n|}{n}<\infty
\]
\cite[Theorem~3.25]{Araujo26}.  The same conclusion holds for the strict
convention under this hypothesis.  Indeed, the two survivor sets differ
only on a subset of
\[
 D:=\{n\in A:0\in X_n\},
\]
and \(\sum_{n\in D}1/n<\infty\).  To see directly that \(D\) has natural
density zero, fix \(M<N\) and note that
\[
 \frac{|D\cap[1,N]|}{N}
 \le\frac{|D\cap[1,M]|}{N}
   +\sum_{\substack{n\in D\\n>M}}\frac1n;
\]
first let \(N\to\infty\), and then \(M\to\infty\).
Our construction necessarily lies outside that summable regime: every
installed scale contributes a fixed positive amount to the sum; see
Remark~\ref{rem:nonsummable}.

The proof has two parts.  Section~\ref{sec:block} constructs a finite block
at scale \(Q=2^j\).  It deletes at least \(3Q/8\) integers from
\([11Q/10,19Q/10]\), even though the periodic union of the responsible
residue classes has measure \(e^{-\Omega(\sqrt j)}\).
Section~\ref{sec:global} assembles long runs of these blocks.  Between
consecutive runs we leave a gap long enough for the finite past periodic
system to recover.  Later moduli lie beyond each earlier recovery cutoff,
so the activation condition prevents them from changing the earlier
averages.

\section{Periodic sets and profinite notation}\label{sec:periodic}
Let \(\Zh\) be the profinite completion of \(\Z\), equipped with Haar
probability measure \(\mu\), and let
\(\pi_q:\Zh\to\Z/q\Z\) be the canonical projection.  For
\(\bar a\in\Z/q\Z\), write
\[
 [\bar a]_q:=\pi_q^{-1}(\bar a).
\]
If \(a\in\Z\), then \([a]_q\) means \([a\bmod q]_q\).  Thus
\(\mu([\bar a]_q)=1/q\).  The trace on \(\Z\) of a finite union of such
cylinders is periodic, and its natural density equals its Haar measure.
We shall also use, for every \(N\ge1\), the elementary bound
\[
 H_N:=\sum_{m=1}^N\frac1m\le1+\log N.
\]

\begin{lemma}[Periodic recovery]\label{lem:periodic}
Let \(\mathcal F\) be a finite family of pairs \((q,Y_q)\), with
\(Y_q\subseteq\Z/q\Z\), and set
\[
 U_{\mathcal F}
 :=\bigcup_{(q,Y_q)\in\mathcal F}\ \bigcup_{a\in Y_q}[a]_q
\]
and
\[
 B_{\mathcal F}
 :=\{m\in\N:m\bmod q\notin Y_q
       \text{ whenever }(q,Y_q)\in\mathcal F\text{ and }q<m\}.
\]
Then
\[
 \sum_{\substack{m<x\\m\in B_{\mathcal F}}}\frac1m
 =\bigl(1-\mu(U_{\mathcal F})\bigr)\log x+O_{\mathcal F}(1)
 \qquad(x\to\infty).
\]
\end{lemma}

\begin{proof}
If \(\mathcal F\) is empty, the assertion is the standard estimate
\(H_{\lceil x\rceil-1}=\log x+O(1)\).  Suppose henceforth that
\(\mathcal F\) is nonempty.  Let
\[
 q_{\max}=\max_{(q,Y_q)\in\mathcal F}q,
 \qquad
 L=\operatorname{lcm}\{q:(q,Y_q)\in\mathcal F\}.
\]
For every \(m>q_{\max}\), all the moduli in \(\mathcal F\) are active.
Therefore \(B_{\mathcal F}\) differs in only finitely many places from the
integer trace of \(\Zh\setminus U_{\mathcal F}\).

That trace is a union of, say, \(r\) residue classes modulo \(L\).  Its
Haar measure is \(r/L=1-\mu(U_{\mathcal F})\).  For a fixed representative
\(a\in\{1,\ldots,L\}\), with \(a=L\) representing the zero class,
\[
 \sum_{\substack{m<x\\m\equiv a\;(\mathrm{mod}\,L)}}\frac1m
 =\sum_{\substack{\ell\ge0\\a+\ell L<x}}\frac1{a+\ell L}
 =\frac1L\log x+O_L(1).
\]
The last estimate follows, for example, by comparing the decreasing
function \(t\mapsto(a+Lt)^{-1}\) with its integral.  Summing over the
\(r\) occupied residue classes and restoring the finite initial
discrepancy proves the lemma.
\end{proof}

\section{A finite dyadic deletion block}\label{sec:block}

For an even positive integer \(k\), define the central family
\[
\cS_k:=\left\{S\subseteq\{1,\ldots,k\}:
       \left||S|-\frac{k}{2}\right|\le\sqrt{k}\right\}.
\]

\begin{lemma}[Arithmetic skeleton]\label{lem:skeleton}
For all sufficiently large integers \(j\), put
\[
 Q=2^j,\qquad k=2\left\lfloor\frac{\sqrt j}{8}\right\rfloor .
\]
There are distinct primes \(p_1,\ldots,p_k\) and, for every
\(S\in\cS_k\), a positive integer \(q_S\) such that
\begin{equation}\label{eq:skeleton-properties}
 \frac{19Q}{20}\le q_S\le\frac{21Q}{20},
 \qquad
 p_i\mid q_S\ \Longleftrightarrow\ i\in S .
\end{equation}
The integers \(q_S\) are pairwise distinct.  Moreover, if
\[
 J=[11Q/10,19Q/10]\cap\Z ,
\]
then \(J\subset(q_S,2q_S]\) and the diameter of \(J\) is smaller than
\(q_S\), for every \(S\in\cS_k\).
\end{lemma}

\begin{proof}
Bertrand's postulate \cite{HardyWright} supplies a prime \(p_i\) satisfying
\[
 4^{k+i}<p_i<2\cdot4^{k+i}\qquad(1\le i\le k).
\]
These intervals are pairwise disjoint, so the primes are distinct.  Set
\(P=\prod_{i=1}^k p_i\).  Since \(p_i<2^{2(k+i)+1}\),
\[
 \log_2 P
 <\sum_{i=1}^k\bigl(2(k+i)+1\bigr)
 =3k^2+2k.
\]
The choice of \(k\) gives \(k\le\sqrt j/4\), and hence
\[
 \frac{P^2}{Q}
 \le 2^{6k^2+4k-j}
 \le 2^{-5j/8+\sqrt j}\longrightarrow0
 \qquad(j\to\infty).
\]

For \(S\in\cS_k\), let \(d_S=\prod_{i\in S}p_i\).  Choose \(R_S\) to be
a closest integer to \(Q/d_S\) in the progression \(1+P\Z\), and put
\(q_S=d_SR_S\).  This choice satisfies
\[
 \left|R_S-\frac{Q}{d_S}\right|\le\frac P2,
 \qquad
 |q_S-Q|\le\frac{d_SP}{2}\le\frac{P^2}{2}.
\]
It remains to check that \(R_S\) is positive.  Since \(d_S\le P\) and
\(Q/P^2\to\infty\), for large \(j\) we have
\[
 \frac{Q}{d_S}\ge\frac QP>P.
\]
Thus \(R_S\ge Q/d_S-P/2>P/2>0\).  The estimate
\(P^2/Q\to0\) now gives the first pair of inequalities in
\eqref{eq:skeleton-properties}.

If \(i\in S\), then \(p_i\mid d_S\), so \(p_i\mid q_S\).  If
\(i\notin S\), then \(p_i\nmid d_S\) and
\(R_S\equiv1\pmod {p_i}\), so \(p_i\nmid q_S\).  This proves the second
part of \eqref{eq:skeleton-properties}; it also shows that distinct sets
\(S\) give distinct integers \(q_S\).

Finally,
\[
 q_S\le\frac{21Q}{20}<\frac{11Q}{10}
 \quad\text{and}\quad
 2q_S\ge\frac{19Q}{10}.
\]
Consequently \(J\subset(q_S,2q_S]\).  Its diameter is at most \(4Q/5\),
which is smaller than \(q_S\ge19Q/20\).
\end{proof}

Fix \(j\) large enough for Lemma~\ref{lem:skeleton}, and retain its
notation.  For every \(1\le i\le k\) and
\(b\in\Z/p_i\Z\), choose an independent Bernoulli random variable
\(\varepsilon_i(b)\) with parameter \(1/2\).  For \(m\in\Z\), set
\[
 K(m):=\{i:\varepsilon_i(m\bmod p_i)=1\}.
\]
Let \(\Pp\) and \(\E\) denote probability and expectation with respect to
this finite family of random labels.
Define the random endpoint set and its periodic footprint by
\[
 E:=\{m\in J:K(m)\in\cS_k\},\qquad
 q_m:=q_{K(m)}\quad(m\in E),\qquad
 U:=\bigcup_{m\in E}[m]_{q_m}.
\]

\begin{lemma}[Endpoint abundance]\label{lem:endpoints}
For all sufficiently large \(j\),
\[
 \Pp\left(|E|<\frac{|J|}{2}\right)
 \le \exp\left(-\frac{4^k}{8k}\right).
\]
\end{lemma}

\begin{proof}
For a fixed \(m\), the variable \(|K(m)|\) has distribution
\(\operatorname{Bin}(k,1/2)\), with mean \(k/2\) and variance \(k/4\).
Chebyshev's inequality therefore gives
\[
 \Pp\bigl(K(m)\notin\cS_k\bigr)
 =\Pp\left(\left||K(m)|-\frac k2\right|>\sqrt k\right)
 \le\frac14.
\]
For \(Z=|E|/|J|\), averaging this estimate over \(m\in J\) gives
\(\E Z\ge3/4\).

Changing one variable \(\varepsilon_i(b)\) affects only those \(m\in J\)
with \(m\equiv b\pmod {p_i}\).  There are at most
\(|J|/p_i+1\) such integers.  The proof of Lemma~\ref{lem:skeleton} gives
\(p_i\le P=o(Q^{1/2})\), whereas \(|J|\asymp Q\); hence
\(|J|\ge p_i\) for large \(j\).  The corresponding bounded-difference
constant for \(Z\) is therefore at most \(2/p_i\).  For each \(i\), there
are exactly \(p_i\) variables \(\varepsilon_i(b)\).  Hence the sum of the
squared bounded-difference constants is at most
\[
 \sum_{i=1}^k p_i\left(\frac2{p_i}\right)^2
 =4\sum_{i=1}^k\frac1{p_i}
 <\frac{k}{4^k},
\]
because every \(p_i>4^{k+1}\).  McDiarmid's bounded-differences inequality
\cite{McDiarmid89}, applied with downward deviation \(1/4\), now yields
\[
 \Pp(Z<1/2)
 \le \exp\left(
       -\frac{2(1/4)^2}{\,k/4^k\,}\right)
 =\exp\left(-\frac{4^k}{8k}\right).
\]
All the independent variables here form a finite family, and the displayed
sum verifies the bounded-difference hypothesis.
\end{proof}

\begin{lemma}[Small periodic footprint]\label{lem:footprint}
For all sufficiently large \(j\),
\[
 \Pp\bigl(\mu(U)>e^{-k/100}\bigr)\le3e^{-k/100}.
\]
\end{lemma}

\begin{proof}
Fix \(\omega\in\Zh\).  For \(S\in\cS_k\), let
\(a_S(\omega)\in\{1,\ldots,q_S\}\) represent \(\omega\bmod q_S\), with
the zero residue represented by \(q_S\), and set
\[
 m_S(\omega):=q_S+a_S(\omega)\in(q_S,2q_S].
\]
Because \(J\subset(q_S,2q_S]\) and has diameter smaller than \(q_S\),
each residue class modulo \(q_S\) has at most one representative in \(J\).
It follows directly from the definitions that
\begin{equation}\label{eq:candidate-equivalence}
 \omega\in U
 \quad\Longleftrightarrow\quad
 \text{there is an \(S\in\cS_k\) for which }
 m_S(\omega)\in J\ \text{and}\ K(m_S(\omega))=S.
\end{equation}

For \(S\in\cS_k\) and \(i\notin S\), define the collision set
\[
 C_{S,i}:=\{\omega\in\Zh:
       m_S(\omega)\equiv\omega\pmod {p_i}\}.
\]
By \eqref{eq:skeleton-properties}, \((q_S,p_i)=1\).  Conditional on any
residue modulo \(q_S\), the displayed congruence prescribes exactly one
residue modulo \(p_i\).  The Chinese remainder theorem, applied to the
uniform residue modulo \(q_Sp_i\), therefore gives
\(\mu(C_{S,i})=1/p_i\).  The set
\[
 C:=\bigcup_{\substack{S\in\cS_k\\i\notin S}}C_{S,i},
\]
depends only on the arithmetic skeleton and is therefore independent of all
the random labels.  The union bound and \(p_i>4^{k+1}\) give
\begin{equation}\label{eq:collision-bound}
 \mu(C)
 \le|\cS_k|\sum_{i=1}^k\frac1{p_i}
 \le2^k\frac{k}{4^{k+1}}
 =\frac{k}{2^{k+2}}.
\end{equation}

Now fix \(\omega\notin C\).  Condition on the \(k\) anchor bits
\(\varepsilon_i(\omega\bmod p_i)\), and write
\[
 T:=\{i:\varepsilon_i(\omega\bmod p_i)=1\}.
\]
If \(i\in S\), then \(p_i\mid q_S\), and hence
\(m_S(\omega)\equiv\omega\pmod {p_i}\).  Therefore
\(K(m_S(\omega))=S\) forces \(S\subseteq T\).  If \(i\notin S\), the
assumption \(\omega\notin C_{S,i}\) says that the residue queried at
\(m_S(\omega)\bmod p_i\) differs from the anchor residue
\(\omega\bmod p_i\).  For a fixed \(S\), these \(k-|S|\) queried variables
have different first coordinates \(i\); each is also distinct from every
conditioned anchor variable.  They are therefore mutually independent and
remain unbiased after the conditioning.  Consequently, conditional
on the anchor bits, the probability that \(S\) contributes to
\eqref{eq:candidate-equivalence} is zero unless \(S\subseteq T\), and in
all cases it is at most
\begin{equation}\label{eq:one-candidate}
 2^{-(k-|S|)}.
\end{equation}
The requirement \(m_S(\omega)\in J\) can only decrease this probability.

Unconditionally, \(|T|\sim\operatorname{Bin}(k,1/2)\).  Hoeffding's
inequality \cite{Hoeffding63} gives
\begin{equation}\label{eq:T-tail}
 \Pp(|T|>3k/5)
 \le\exp\left(-\frac{2(k/10)^2}{k}\right)
 =e^{-k/50}.
\end{equation}
Suppose that \(|T|\le3k/5\) and that \(S\in\cS_k\) is contained in \(T\).
The map \(S\mapsto T\setminus S\) is injective, and
\[
 |T\setminus S|
 \le\frac{3k}{5}-\left(\frac k2-\sqrt k\right)
 =\frac{k}{10}+\sqrt k.
\]
Let \(N_k=\lfloor3k/5\rfloor\) and
\(M_k=\lfloor k/10+\sqrt k\rfloor\).  Padding \(T\) to a set with
\(N_k\) elements is possible because the integer \(|T|\) is at most
\(\lfloor3k/5\rfloor=N_k\).  Since \(S\mapsto T\setminus S\) is
injective, the number of possible candidates is at most
\[
 \sum_{\ell=0}^{M_k}\binom{N_k}{\ell}.
\]

We record the entropy estimate used to bound this sum.  If
\(0<\rho\le1/2\) and \(M=\lfloor\rho N\rfloor\), then
\begin{equation}\label{eq:entropy-bound}
 \sum_{\ell=0}^{M}\binom N\ell
 \le\exp\bigl(NH(\rho)\bigr),
 \qquad
 H(\rho)=-\rho\log\rho-(1-\rho)\log(1-\rho).
\end{equation}
Indeed, for \(\ell\le\rho N\), the quantity
\(\rho^\ell(1-\rho)^{N-\ell}\) is at least
\(\rho^{\rho N}(1-\rho)^{(1-\rho)N}=e^{-NH(\rho)}\).
Multiplication by this lower bound and comparison with the complete
binomial expansion proves \eqref{eq:entropy-bound}.

Here \(M_k/N_k\to1/6\) and \(N_k/k\to3/5\).  In particular,
\(0<M_k/N_k<1/2\) for sufficiently large \(k\), so
\eqref{eq:entropy-bound} applies with \(\rho=M_k/N_k\).  The numerical
margin needed below follows from
\[
 \frac35H(1/6)
 =\frac1{10}\log6+\frac12\log(6/5)
 <\frac1{10}\frac95+\frac12\frac15
 =\frac7{25}=0.28.
\]
The inequality \(\log(6/5)<1/5\) is the standard
\(\log(1+t)<t\) with \(t=1/5\).  Also \(\log6<9/5\), because the first
seven terms of the positive power series for \(e^{9/5}\) have sum greater
than \(6\).  By continuity of \(H\), for all sufficiently large \(k\),
\begin{equation}\label{eq:candidate-count}
 \#\{S\in\cS_k:S\subseteq T\}\le e^{0.29k}.
\end{equation}

For the same range of \(k\), \(\sqrt k\le k/100\), so every
\(S\in\cS_k\) satisfies \(k-|S|\ge0.49k\).  The positive series
\[
 \log2
 =2\sum_{r=0}^{\infty}
   \frac1{(2r+1)3^{2r+1}}
 >\frac{56}{81}
\]
shows that \(0.49\log2>0.33\).  Hence \eqref{eq:one-candidate} is at most
\(e^{-0.33k}\).  For every fixed realization of the anchor bits with
\(|T|\le3k/5\), the conditional union bound, together with
\eqref{eq:one-candidate} and \eqref{eq:candidate-count}, gives
\[
 \Pp\bigl(\omega\in U\mid
   \varepsilon_i(\omega\bmod p_i),\ 1\le i\le k\bigr)
 \le e^{0.29k}e^{-0.33k}=e^{-0.04k}.
\]
Averaging over the anchor bits therefore yields, uniformly for
\(\omega\notin C\),
\[
 \Pp(\omega\in U,\ |T|\le3k/5)\le e^{-0.04k}.
\]

The indicator of \(U\) is nonnegative, so Tonelli's theorem gives
\[
 \E_{\varepsilon}\mu(U)
 =\int_{\Zh}\Pp_{\varepsilon}(\omega\in U)\,d\mu(\omega).
\]
Split the integral over \(C\) and \(\Zh\setminus C\), and use
\eqref{eq:collision-bound}, \eqref{eq:T-tail}, and the preceding estimate.
For all sufficiently large \(k\),
\[
 \E\mu(U)
 \le\frac{k}{2^{k+2}}+e^{-k/50}+e^{-0.04k}
 \le3e^{-k/50}.
\]
Markov's inequality, with threshold \(e^{-k/100}\), now gives
\[
 \Pp\bigl(\mu(U)>e^{-k/100}\bigr)
 \le e^{k/100}\E\mu(U)
 \le3e^{-k/100},
\]
as required.
\end{proof}

\begin{lemma}[Finite block lemma]\label{lem:block}
For every sufficiently large integer \(j\), put
\[
 Q=2^j,\qquad
 k_j=2\left\lfloor\frac{\sqrt j}{8}\right\rfloor,\qquad
 \eta_j=e^{-k_j/100}.
\]
There are an endpoint set \(E_j\subseteq\N\), a modulus \(q_{j,m}\) for
each \(m\in E_j\), and a periodic cylinder union
\[
 U_j=\bigcup_{m\in E_j}[m]_{q_{j,m}}
\]
such that
\begin{equation}\label{eq:block-properties}
\begin{gathered}
 E_j\subseteq[11Q/10,19Q/10]\cap\Z,\qquad |E_j|\ge3Q/8,\\
 \frac{19Q}{20}\le q_{j,m}\le\frac{21Q}{20}<m
 \quad(m\in E_j),\qquad
 \mu(U_j)\le\eta_j .
\end{gathered}
\end{equation}
Moreover, \(\sum_j\eta_j<\infty\).
\end{lemma}

\begin{proof}
Lemmas~\ref{lem:endpoints} and \ref{lem:footprint} show that the probability
of failure of either
\[
 |E|\ge|J|/2,\qquad \mu(U)\le e^{-k/100}
\]
is at most
\[
 \exp\left(-\frac{4^k}{8k}\right)+3e^{-k/100},
\]
which is smaller than \(1\) for all sufficiently large \(k\).  Fix a
labelling for which both conclusions hold, and use it to define
\(E_j,q_{j,m}\), and \(U_j\).  Since \(|J|\ge4Q/5-1\), we have
\(|E_j|\ge3Q/8\) for large \(Q\).  Lemma~\ref{lem:skeleton} gives the
modulus bounds, and the strict inequality \(q_{j,m}<m\) follows from
\[
 q_{j,m}\le21Q/20<11Q/10\le m.
\]
This proves \eqref{eq:block-properties}.

Finally,
\[
 k_j\ge\frac{\sqrt j}{4}-2,
 \qquad
 \eta_j\le e^{1/50}e^{-\sqrt j/400}.
\]
To see that the last majorant is summable, group the integers
\(r^2\le j<(r+1)^2\); the contribution of the \(r\)-th group is at most
\((2r+1)e^{-r/400}\), up to the fixed factor \(e^{1/50}\).  The resulting
series converges.
\end{proof}

\section{The global construction}\label{sec:global}

Set \(\epsilon=1/100\).  Choose \(j_0\) sufficiently large that
Lemma~\ref{lem:block} holds for every \(j\ge j_0\) and
\begin{equation}\label{eq:global-smallness}
 \sum_{j\ge j_0}\eta_j<\epsilon,
 \qquad
 \frac1{(2j_0+1)\log2}<\epsilon.
\end{equation}
For each \(j\ge j_0\), fix once and for all a block
\((E_j,(q_{j,m})_{m\in E_j},U_j)\) supplied by that lemma.

\subsection{Choosing the epochs}

We recursively choose integers \(a_1<a_2<\cdots\) and install the blocks
with indices in
\[
 I_t:=\{a_t,a_t+1,\ldots,2a_t\}
\]
during epoch \(t\).  Suppose that the first \(t-1\) epochs have already
been chosen.  For every modulus \(q\) occurring in those epochs, put
\[
 Y_q^{(t-1)}
 :=\{m\bmod q:s<t,\ j\in I_s,\ m\in E_j,\ q_{j,m}=q\},
\]
and let \(\mathcal F_{t-1}\) be the finite family of all pairs
\((q,Y_q^{(t-1)})\) with \(Y_q^{(t-1)}\ne\varnothing\).  For \(t=1\),
take \(\mathcal F_0=\varnothing\).  Grouping equal moduli changes neither
the survivor set nor the union of cylinders, and hence
\[
 U_{\mathcal F_{t-1}}
 =V_{t-1}:=\bigcup_{\substack{s<t\\j\in I_s}}U_j.
\]
By induction, the already chosen intervals \(I_s\), \(s<t\), are pairwise
disjoint subsets of \(\{j:j\ge j_0\}\).  The union bound and
\eqref{eq:global-smallness} therefore imply
\[
 \mu(V_{t-1})
 \le\sum_{\substack{s<t\\j\in I_s}}\eta_j
 <\epsilon.
\]

By Lemma~\ref{lem:periodic}, the logarithmic averages of
\(B_{\mathcal F_{t-1}}\) converge to
\(1-\mu(V_{t-1})>1-\epsilon\).  We may therefore choose \(a_t\) so large
that
\begin{equation}\label{eq:epoch-choice}
\begin{gathered}
 a_t\ge j_0,\qquad
 a_t>2a_{t-1}+3\quad(t\ge2),\\
 \frac1{\log x_t}
 \sum_{\substack{m<x_t\\m\in B_{\mathcal F_{t-1}}}}\frac1m
 \ge1-2\epsilon=\frac{49}{50},
 \qquad x_t:=2^{a_t-1}.
\end{gathered}
\end{equation}
Indeed, the convergence just noted makes the last inequality true for all
sufficiently large \(x\), so it remains true after restricting \(x\) to
the unbounded sequence of dyadic numbers \(2^{a-1}\).  This completes the
recursion.  In particular, the epochs are pairwise disjoint and
\(a_t\to\infty\).

\subsection{The fixed congruence system}

Let
\[
 \cI:=\bigcup_{t\ge1}I_t
\]
be the set of installed scales.  Modulus ranges at distinct scales are
disjoint, since
\begin{equation}\label{eq:scale-separation}
 \frac{21}{20}2^j<\frac{19}{20}2^{j+1}.
\end{equation}
For each modulus \(q\) appearing in an installed block, define
\[
 X_q:=\{m\bmod q:j\in\cI,\ m\in E_j,\ q_{j,m}=q\}
 \subseteq\Z/q\Z,
\]
and let \(A\) be the set of those moduli.  Repeated occurrences of the same
modulus inside one block are thereby grouped into a single residue set.
Equation~\eqref{eq:scale-separation} shows that a modulus cannot occur at
two different scales.

The family is infinite.  Indeed, \(|E_j|\ge3\cdot2^j/8>0\) at every
installed scale, so every such scale supplies a modulus; by
\eqref{eq:scale-separation}, these moduli are distinct as the scale varies.
Let \(B=B(A,X)\) be the survivor set defined in \eqref{eq:defB}.  Once the
sequence \((a_t)\) has been chosen, both \(A\) and every \(X_q\) are fixed;
they do not depend on the averaging parameter.

\begin{remark}[Failure of the summability hypothesis]\label{rem:nonsummable}
Within a fixed block, two distinct endpoints assigned the same modulus
represent distinct residue classes: their distance is smaller than that
modulus by Lemma~\ref{lem:skeleton}.  Together with
\eqref{eq:scale-separation}, this gives, for every installed \(j\),
\[
 \sum_{\substack{q\in A\\19\cdot2^j/20\le q\le21\cdot2^j/20}}
       \frac{|X_q|}{q}
 =\sum_{m\in E_j}\frac1{q_{j,m}}
 \ge\frac{|E_j|}{21\cdot2^j/20}
 \ge\frac5{14}.
\]
There are infinitely many installed scales, and therefore
\(\sum_{q\in A}|X_q|/q=\infty\).  Thus the construction does not overlap
the summable positive result cited in the introduction.
\end{remark}

\subsection{Recovery cutoffs}

At \(x_t=2^{a_t-1}\), consider a modulus belonging to an epoch \(s\ge t\).
Its scale \(j\) satisfies \(j\ge a_s\ge a_t\), so
\[
 q\ge\frac{19}{20}2^j
 \ge\frac{19}{20}2^{a_t}
 >2^{a_t-1}=x_t.
\]
Such a modulus is not active at any integer \(m<x_t\).  Consequently, below
\(x_t\) the final survivor \(B\) agrees exactly with the finite-past
survivor \(B_{\mathcal F_{t-1}}\): by
\eqref{eq:scale-separation}, no later scale can add a new residue class to
a modulus already present in \(\mathcal F_{t-1}\), and all genuinely future
moduli are inactive below \(x_t\).  By \eqref{eq:epoch-choice},
\[
 L_B(x_t)\ge\frac{49}{50}.
\]
Since \(x_t\to\infty\), it follows that
\begin{equation}\label{eq:limsup}
 \limsup_{x\to\infty}L_B(x)\ge\frac{49}{50}.
\end{equation}

\subsection{Deletion cutoffs}

Put \(y_t=2^{2a_t+1}\).  If \(j\in I_t\) and \(m\in E_j\), then
\[
 m\le\frac{19}{10}2^j
 \le\frac{19}{10}2^{2a_t}
 <2^{2a_t+1}=y_t.
\]
Moreover, Lemma~\ref{lem:block} gives
\[
 q_{j,m}\le\frac{21}{20}2^j
 <\frac{11}{10}2^j\le m.
\]
By the definition of \(X_{q_{j,m}}\), one has
\(m\bmod q_{j,m}\in X_{q_{j,m}}\).  Since \(q_{j,m}<m\), this modulus is
active at \(m\), and consequently \(m\notin B\).  Thus every point of
\(E_j\) is absent from \(B\cap[1,y_t)\).

The endpoint intervals at consecutive scales are disjoint:
\[
 \frac{19}{10}2^j<\frac{11}{10}2^{j+1}.
\]
Hence the deleted endpoints supplied by different \(j\in I_t\) are
distinct.  At a single scale, \eqref{eq:block-properties} gives
\[
 \sum_{m\in E_j}\frac1m
 \ge\frac{|E_j|}{19\cdot2^j/10}
 \ge\frac{3\cdot2^j/8}{19\cdot2^j/10}
 =\frac{15}{76}.
\]
Because \(I_t\) contains \(a_t+1\) scales, we obtain
\begin{equation}\label{eq:deletion-average}
 L_B(y_t)
 \le\frac{H_{y_t-1}}{\log y_t}
 -\frac{15}{76}\,
   \frac{a_t+1}{(2a_t+1)\log2}.
\end{equation}

Since \(a_t\ge j_0\), the harmonic-number bound from
Section~\ref{sec:periodic} and \eqref{eq:global-smallness} show that
\[
 \frac{H_{y_t-1}}{\log y_t}
 \le1+\frac1{\log y_t}<1+\epsilon.
\]
Also \((a_t+1)/(2a_t+1)>1/2\), and
\[
 \frac{15}{152\log2}>\frac18.
\]
For completeness, \(\log2<3/4\), since the first four terms in the
positive series for \(e^{3/4}\) already sum to more than \(2\); hence the
left side is greater than \(15/114=5/38>1/8\).  Substitution into
\eqref{eq:deletion-average} gives
\[
 L_B(y_t)
 <1+\frac1{100}-\frac18
 =\frac{177}{200}.
\]
Since \(y_t\to\infty\),
\begin{equation}\label{eq:liminf}
 \liminf_{x\to\infty}L_B(x)\le\frac{177}{200}.
\end{equation}

Equations~\eqref{eq:limsup} and \eqref{eq:liminf} prove
Theorem~\ref{thm:main}.

\section{Concluding remarks}

The construction separates two effects that coincide for many classical
sieves.  Within an installed scale, the active tails delete a fixed amount
of harmonic mass.  Globally, however, the completed periodic footprint of
that scale has summably small Haar measure.  Long deletion epochs exploit
the first fact, while the recovery cutoffs exploit the second.

Erd\H{o}s Problem 25 is the singleton special case of Problem 486, in
which \(|X_q|=1\) for every \(q\in A\) \cite{Bloom25}.  The present
argument does not settle that problem.  Indeed, the finite block in
Section~\ref{sec:block} groups many endpoint residues into the same set
\(X_q\); this multi-residue feature makes it possible to obtain a fixed
local harmonic deletion while keeping the completed periodic footprint
small.

\begin{thebibliography}{99}

\bibitem{Araujo26}
F.~Ara\'ujo,
\newblock Sarnak's program for Erd\H{o}s sieves. Part I:
topological dynamics and light tails,
\newblock preprint, 2026,
\newblock \href{https://arxiv.org/abs/2602.24031}
{arXiv:2602.24031}.

\bibitem{Be34}
A.~S. Besicovitch,
\newblock On the density of certain sequences of integers,
\newblock \emph{Math. Ann.} \textbf{110} (1934), 336--341,
\newblock \href{https://doi.org/10.1007/BF01448032}
{doi:10.1007/BF01448032}.

\bibitem{Be35}
F.~Behrend,
\newblock On sequences of numbers not divisible one by another,
\newblock \emph{J. London Math. Soc.} \textbf{10} (1935), 42--44,
\newblock \href{https://doi.org/10.1112/jlms/s1-10.37.42}
{doi:10.1112/jlms/s1-10.37.42}.

\bibitem{Bloom25}
T.~F. Bloom,
\newblock Erd\H{o}s Problem 25,
\newblock \url{https://www.erdosproblems.com/25},
accessed July 15, 2026.

\bibitem{DaEr36}
H.~Davenport and P.~Erd\H{o}s,
\newblock On sequences of positive integers,
\newblock \emph{Acta Arith.} \textbf{2} (1936), 147--151,
\newblock \href{https://doi.org/10.4064/aa-2-1-147-151}
{doi:10.4064/aa-2-1-147-151}.

\bibitem{DaEr51}
H.~Davenport and P.~Erd\H{o}s,
\newblock On sequences of positive integers,
\newblock \emph{J. Indian Math. Soc. (N.S.)} \textbf{15} (1951), 19--24,
\newblock \href{https://doi.org/10.18311/JIMS/1951/17063}
{doi:10.18311/JIMS/1951/17063}.

\bibitem{Er57}
P.~Erd\H{o}s,
\newblock On the distribution function of additive arithmetical functions
and on some related problems,
\newblock \emph{Rend. Sem. Mat. Fis. Milano} \textbf{27} (1957), 45--49,
\newblock \href{https://doi.org/10.1007/BF02922565}
{doi:10.1007/BF02922565}.

\bibitem{Er61}
P.~Erd\H{o}s,
\newblock Some unsolved problems,
\newblock \emph{Publ. Math. Inst. Hung. Acad. Sci., Ser. A}
\textbf{6} (1961), 221--254.

\bibitem{HardyWright}
G.~H. Hardy and E.~M. Wright,
\newblock \emph{An Introduction to the Theory of Numbers},
\newblock sixth ed., revised by D.~R. Heath-Brown and J.~H. Silverman,
Oxford University Press, Oxford, 2008.

\bibitem{Hoeffding63}
W.~Hoeffding,
\newblock Probability inequalities for sums of bounded random variables,
\newblock \emph{J. Amer. Statist. Assoc.} \textbf{58} (1963), no.~301,
13--30,
\newblock \href{https://doi.org/10.1080/01621459.1963.10500830}
{doi:10.1080/01621459.1963.10500830}.

\bibitem{McDiarmid89}
C.~McDiarmid,
\newblock On the method of bounded differences,
\newblock in \emph{Surveys in Combinatorics, 1989},
London Math. Soc. Lecture Note Ser. \textbf{141},
Cambridge University Press, Cambridge, 1989, pp.~148--188,
\newblock \href{https://doi.org/10.1017/CBO9781107359949.008}
{doi:10.1017/CBO9781107359949.008}.

\end{thebibliography}

\end{document}
